
\renewcommand\thesection{\arabic{section}}
\renewcommand{\contentsname}{Cuprins}
\renewcommand{\figurename}{Figura}
\renewcommand{\tablename}{Tabel}

\setcounter{tocdepth}{6}
\setcounter{secnumdepth}{6}


% \usepackage[backend=bibtex,style=numeric,citestyle=numeric,backref=true,sorting=none]{biblatex}




\titleformat{\section}{\normalfont\Large\justifyheading}{\thesection}{18pt}{}



\definecolor{codegreen}{rgb}{0,0.6,0}
\definecolor{codegray}{rgb}{0.5,0.5,0.5}
\definecolor{codepurple}{rgb}{0.58,0,0.82}
\definecolor{backcolour}{rgb}{0.96,0.96,0.96}

\lstdefinestyle{mystyle}{
	backgroundcolor=\color{backcolour},   
	commentstyle=\color{codegreen},
	keywordstyle=\color{blue},
	numberstyle=\tiny\color{codegray},
	stringstyle=\color{codepurple},
	basicstyle=\fontsize{10}{10}\selectfont\ttfamily,
	breakatwhitespace=false,         
	breaklines=true,                 
	captionpos=b,                    
	keepspaces=true,                 
	numbers=left,                    
	numbersep=2pt,                  
	showspaces=false,                
	showstringspaces=false,
	showtabs=false,                  
	tabsize=2
}

\lstset{style=mystyle}




\setmainfont{UTSans-Regular.otf}





\newpage



\newpage
\pagenumbering{arabic}
\renewcommand{\bibname}{Bibliografie}
\renewcommand\lstlistlistingname{\normalsize{CODURI SURSĂ}}
\renewcommand{\lstlistingname}{\normalsize{Secvență de Cod}}
\renewcommand{\listfigurename}{FIGURI}
\renewcommand{\cftloftitlefont}{\normalsize}
\renewcommand{\cftfigfont}{\normalsize}
\renewcommand{\cftfigpagefont}{\normalsize}
\renewcommand{\listtablename}{TABELE}
\renewcommand{\cftlottitlefont}{\normalsize}
\renewcommand{\cfttabfont}{\normalsize}
\renewcommand{\cfttabpagefont}{\normalsize}
\clearpage

% % trebuie completata manual





\setcounter{page}{18}

\begin{comment}
	
	
	\section{Pașii de implementări}
	Pentru a putea crea un geamăn digital pentru brațul robotic a trebuit realizate următorii pași:
	
	\begin{itemize}
		\item Pentru primul pas a fost necesar crearea modelului 3D al brațului
		\item Al doilea pas a fost exportarea modelului 3D din Fusion în format de fișier URDF și importare în Unity
		\item Al treilea pas după impotare a fost setarea valorilor de amortizare, frecare, și limitele superioare și inferioare.
		\item Al patrulea pas a constat în aplicarea de constrângeri între componente în jurul motoarelor.
		\item Al cincilea pas după constrângeri a fost crearea sistemului de control
		\item Al saselea pas a constat în transmiterea datelor provenite de la robotul din Unity la robotul real
		\item Al saptelea pas a constat în aplicarea valorilor primite asupra motoarelor și realizarea de modificări ulterioare.
	\end{itemize}
\end{comment}

\section(Arhitectura sistemului)







